\makesectionpage{Deep Learning}


\makesubsectionpage{Deep Neural Networks}

\begin{frame}{Deep Learning}
\begin{columns}
\begin{column}{0.60\textwidth}
\begin{itemize}
\item Deep Learning ist ein \textbf{Sonderfall des überwachten Lernens}.
\item Die Modelle basieren auf Architekturen \textbf{komplexer neuronaler Netze}.
\item \textbf{deep} bezieht sich auf die Anzahl der \textbf{verborgenen Schichten}. Modelle können Hunderte oder Tausende davon haben.
\item Sie werden mit großen Mengen an \textbf{gekennzeichneten Daten} trainiert.
\item Sie lernen Merkmale oft \textbf{direkt aus den Daten}, ohne manuelle Merkmalsextraktion.
\end{itemize}
\end{column}
\begin{column}{0.35\textwidth}
\resizebox{1.0\linewidth}{!}{%
\centering
\begin{tikzpicture}[every node/.style={font=\footnotesize}]

% KI
\filldraw[fill=PrimaryColor, fill opacity=0.1, draw=PrimaryColor, line width=1pt] (0,0) circle (3.75);
\node[align=center] at (0,3.1)
{\textbf{Künstliche}\\\textbf{Intelligenz}};

% Machine Learning
\filldraw[fill=AccentColor, fill opacity=0.1, draw=AccentColor, line width=1pt] (0,-0.6) circle (3);
\node[align=center] at (0,1.7)
{\textbf{Machine}\\\textbf{Learning}};

% Deep Learning
\filldraw[fill=AlertColor, fill opacity=0.1, draw=AlertColor, line width=1pt] (0,-1.2) circle (2.2);
\node[align=center] at (0,0.4)
{\textbf{Deep}\\\textbf{Learning}};

% Generative AI
\filldraw[fill=ThirdColor, fill opacity=0.1, draw=ThirdColor, line width=1pt] (0,-1.7) circle (1.5);
\node[align=center] at (0,-1.1)
{\textbf{Generative}\\\textbf{KI}};

% LLM
\filldraw[fill=ExampleColor, fill opacity=0.1, draw=ExampleColor, line width=1pt] (0,-2.3) circle (0.75);
\node[align=center] at (0,-2.3)
{\textbf{LLM}};

\end{tikzpicture}%
}
\end{column}
\end{columns}
\source{\cite{mathworks_deep_learning}}
\end{frame}

\begin{frame}{Deep Learning: Rechenleistung}
\begin{columns}
\begin{column}{0.60\textwidth}
Beim Deep Learning werden mit sehr großen Datensätzen \textbf{enorm viele Parameter} neuronaler Netze optimiert.
\begin{itemize}
\item Dadurch erfordert Deep Learning eine \textbf{enorme Rechenleistung}.
\item Diese ist erst seit den \textbf{2000er-Jahren} allgemein verfügbar.
\item \textbf{Leistungsstarke Grafikkarten} bieten eine parallele Architektur, die besonders effizient ist.
\item In der \textbf{Cloud} wird Rechenleistung frei oder gegen Gebühr bereitgestellt.
\end{itemize}
\end{column}
\begin{column}{0.35\textwidth}
\centering
\includesvg[width=0.9\linewidth]{figures/gpu.svg}
\end{column}
\end{columns}
\source{\cite{mathworks_deep_learning}}
\end{frame}


\begin{frame}{Tensoren}
\begin{columns}
\begin{column}{0.60\textwidth}
Ein \textbf{Tensor} ist ein mehrdimensionales Zahlenfeld und die zentrale Datenstruktur im Deep Learning.
\begin{itemize}
\item \textbf{Rang 0}: Skalar, eine einzelne Zahl.
\item \textbf{Rang 1}: Vektor, eine Zahlenreihe.
\item \textbf{Rang 2}: Matrix, eine Tabelle.
\item \textbf{Rang $n$}: Verallgemeinerung auf $n$ Achsen.
\end{itemize}
Die \textbf{Form} (shape) gibt die Größe entlang jeder Achse an.
\end{column}
\begin{column}{0.35\textwidth}
\resizebox{\linewidth}{!}{%
\begin{tikzpicture}[every node/.style={font=\scriptsize}]

% ===== Obere Reihe (vertikal zentriert um y = 3.8) =====

% Skalar (0.5 hoch -> 3.55 bis 4.05)
\filldraw[fill=PrimaryColor!20, draw=PrimaryColor, very thick] (0,3.55) rectangle (0.5,4.05);
\node at (0.25,2.7) {\textbf{Skalar}};
\node at (0.25,2.35) {Rang 0};

% Vektor (1.5 hoch -> 3.05 bis 4.55)
\foreach \y in {3.05,3.55,4.05} {
  \filldraw[fill=AccentColor!20, draw=AccentColor, very thick] (2.3,\y) rectangle (2.8,\y+0.5);
}
\node at (2.55,2.7) {\textbf{Vektor}};
\node at (2.55,2.35) {Rang 1};

% ===== Untere Reihe (vertikal zentriert um y = 0.75) =====

% Matrix (1.5 hoch -> 0 bis 1.5)
\foreach \x in {0,0.5} {
  \foreach \y in {0,0.5,1} {
    \filldraw[fill=AlertColor!20, draw=AlertColor, very thick] (\x,\y) rectangle (0.5+\x,0.5+\y);
  }
}
\node at (0.5,-0.4) {\textbf{Matrix}};
\node at (0.5,-0.75) {Rang 2};

% Tensor (Rang 3, visuell 1.36 hoch -> Basis bei 0.07)
\foreach \z in {0,0.18,0.36} {
  \foreach \x in {0,0.5} {
    \foreach \y in {0,0.5} {
      \filldraw[fill=ThirdColor!20, draw=ThirdColor, very thick]
        (2.3+\x+\z,0.07+\y+\z) rectangle (2.8+\x+\z,0.57+\y+\z);
    }
  }
}
\node at (2.7,-0.4) {\textbf{Tensor}};
\node at (2.7,-0.75) {Rang 3};

\end{tikzpicture}%
}
\end{column}
\end{columns}
\end{frame}


\begin{frame}{Tensoren: Beispiel RGB-Bild}
\begin{columns}
\begin{column}{0.47\textwidth}
Ein \textbf{RGB-Bild} ist ein typischer \textbf{Rang-3-Tensor}:
\begin{itemize}
\item \textbf{Höhe}: Anzahl der Bildzeilen,
\item \textbf{Breite}: Anzahl der Bildspalten,
\item \textbf{Kanäle}: Rot, Grün, Blau.
\end{itemize}
Die Form lautet also \textbf{(Höhe, Breite, 3)}. Jeder Eintrag ist ein Helligkeitswert eines Kanals an einer Position.
\end{column}
\begin{column}{0.48\textwidth}
\resizebox{\linewidth}{!}{%
\begin{tikzpicture}[every node/.style={font=\scriptsize}]

% Drei gestapelte Kanäle
\foreach \c/\name/\dx/\dy in {AlertColor/Rot/0/0, ExampleColor/Grün/0.6/0.5, PrimaryColor/Blau/1.2/1.0} {
  \begin{scope}[shift={(\dx,\dy)}]
    \draw[fill=\c!15, draw=\c, very thick] (0,0) rectangle (3,3);
    \foreach \i in {0.5,1,1.5,2,2.5} {
      \draw[\c!40] (\i,0) -- (\i,3);
      \draw[\c!40] (0,\i) -- (3,\i);
    }
  \end{scope}
}

\node[text=PrimaryColor] at (2.6,1.3) {\textbf{Blau}};
\node[text=ExampleColor] at (2.2,0.8) {\textbf{Grün}};
\node[text=AlertColor] at (1.6,0.3) {\textbf{Rot}};

\draw[<->, thick] (-0.3,0) -- (-0.3,3) node[midway,left,rotate=90,anchor=south] {Höhe};
\draw[<->, thick] (0,-0.3) -- (3,-0.3) node[midway,below] {Breite};

\end{tikzpicture}%
}
\end{column}
\end{columns}
\end{frame}

\begin{frame}{Tensoren im neuronalen Netz}
\begin{columns}
\begin{column}{0.47\textwidth}
\textbf{Datenrepräsentation:}
\begin{itemize}
\item Tensoren stellen Datentypen wie \textbf{Bilder} und \textbf{Texte} dar.
\item Eingabedaten werden in strukturierten Tensoren organisiert.
\end{itemize}

\vspace{0.4em}
\textbf{Gewichte der Schichten:}
\begin{itemize}
\item Die \textbf{Gewichte} neuronaler Netze sind ebenfalls Tensoren.
\item Sie werden während des \textbf{Trainings} angepasst.
\end{itemize}
\end{column}
\begin{column}{0.48\textwidth}
\textbf{Mathematische Operationen:}
\begin{itemize}
\item \textbf{Matrixmultiplikation} von Eingaben und Gewichten,
\item \textbf{Aktivierungsfunktionen} für Nichtlinearitäten,
\item \textbf{Gradientenberechnung} für die Backpropagation.
\end{itemize}

\vspace{0.4em}
\textbf{Vorhersage:}
\begin{itemize}
\item Die Ergebnisse führen zu \textbf{Modellausgaben} und \textbf{Vorhersagen}.
\end{itemize}
\end{column}
\end{columns}
\end{frame}

\begin{frame}{TensorFlow}
\begin{columns}
\begin{column}{0.35\textwidth}
\centering
\includesvg[width=0.9\linewidth]{figures/Tensorflow_logo.svg}
\end{column}
\begin{column}{0.60\textwidth}
\textbf{TensorFlow} ist ein quelloffenes Framework für maschinelles Lernen, entwickelt von Google.
\begin{itemize}
\item Rechnungen werden als \textbf{Datenflussgraphen} über Tensoren beschrieben.
\item Es nutzt \textbf{CPUs, GPUs} und \textbf{TPUs} zur parallelen Verarbeitung.
\item Man hat \textbf{detaillierte Kontrolle} über alle Parameter, etwa zur Entwicklung \textbf{neuartiger Netztopologien}.
\item Wird oft in der \textbf{Forschung} eingesetzt.
\end{itemize}
\end{column}
\end{columns}
\source{\cite{tensorflow2015}}
\end{frame}


\begin{frame}{Keras}
\begin{columns}
\begin{column}{0.60\textwidth}
\textbf{Keras} ist leicht erlernbare API für neuronale Netze, die auf Frameworks wie TensorFlow aufsetzt.
\begin{itemize}
\item Modelle werden aus \textbf{Schichten} einfach zusammengesetzt.
\item Der Fokus liegt auf \textbf{Benutzerfreundlichkeit} und schneller Umsetzung.
\item Für reine \textbf{Anwender*innen} einfacher zu erlernen als TensorFlow.
\item Bietet bereits \textbf{vortrainierte Modelle} für Klassifikationsaufgaben.
\item Ideal für \textbf{Einstieg} und schnelles \textbf{Prototyping}.
\end{itemize}
\end{column}
\begin{column}{0.35\textwidth}
\centering
\includesvg[width=0.9\linewidth]{figures/keras_logo.svg}
\end{column}

\end{columns}
\source{\cite{chollet2015keras}}
\end{frame}

\makequizpage{\QUIZURL}{\QUIZQRCODE}{\QUIZIMAGE}
\makequestionspage{\FRAGENIMAGE}